\documentclass[10pt,a4paper]{amsart}
\usepackage[margin=19mm]{geometry}
\usepackage{amsmath,amssymb,mathtools}
\usepackage[scr=rsfs,cal=euler]{mathalfa}
\usepackage{xfrac}
\usepackage[unicode,hidelinks,bookmarks=false]{hyperref}
\newcommand{\ie}{i.e.\ }
\newcommand{\cf}{cf.\ }
\newcommand{\resp}{respectively\ }
\newcommand{\Subsection}{\subsection}
\providecommand{\nobreakdash}{\nobreak\hbox{-}}
\setlength{\emergencystretch}{2em}
\multlinegap0pt
\usepackage{bm}
\usepackage{bbm}
\usepackage{dsfont}
\usepackage{xspace}
\usepackage{cancel}
\usepackage{mathtools}
\usepackage{amsthm}
\makeatletter
\@ifundefined{theorem}{\newtheorem{theorem}{Theorem}}{}
\makeatother
\title[A Symmetric Polynomial Approach to Poncelet Triangles Inscri]{A Symmetric Polynomial Approach to Poncelet Triangles Inscribed in a circle and circumscribed about Central Conics}
\author{Mohammad Hassan Murad}
\date{Source: arXiv:2603.28822v2 (CC BY 4.0)}
\begin{document}
\maketitle

\noindent\emph{Source and licence.} A Symmetric Polynomial Approach to Poncelet Triangles Inscribed in a circle and circumscribed about Central Conics. Version arXiv:2603.28822v2 (CC BY 4.0), distributed under the Creative Commons Attribution 4.0 International licence, as recorded in the arXiv licence metadata for this exact version and shown on the abstract page https://arxiv.org/abs/2603.28822v2. Full rights notice: licenses/arxiv-2603.28822-CC-BY-4.0.txt.\par\smallskip

\noindent\textit{Editorial scope.} Counted unit: the Theorem whose source label is \texttt{thm:midptdist} in the article (source0.tex lines 1411--1429, source label \texttt{thm:midptdist}), delivered together with the article's own complete original proof of it (source0.tex lines 1431--1531, one \texttt{proof} environment of 101 lines). No mathematics is written, repaired or re-derived: statement and proof are the article's own text line for line. Required context retained uncounted: eq:zHrelation (lines 165--171). Every definition, display and label of those lines is retained unchanged. Omitted from this package and not redistributed: 1--164, 172--1410, 1430--1430, 1532--2150. Nothing inside a retained excerpt is omitted or altered: every retained body is delivered exactly as the article's lines are written, so no omission receipt of any kind accompanies this package. Citations: the retained excerpts carry no citation command at all, so no bibliography entry of the article is transcribed and no reference is redistributed. Reference adaptation: none was needed and none was made; every reference inside the delivered unit resolves inside it, so no unresolved reference or citation is produced. Printed slips kept literal: no printed word, symbol, hypothesis or number of the counted statement or of its proof was altered. Layout and class: the article's own class is \texttt{amsart} with packages including amsaddr, graphicx, pgfplots, orcidlink, geometry, stix, subcaption, caption, amsmath, amsthm, cite, hyperref, url; the wrapper is the lane's pinned \texttt{amsart} editorial wrapper at 10pt on A4, which restates the theorem-like environments the retained text needs and adds the article's own macro definitions that the retained text needs (none), with extra wrapper packages (bm, bbm, dsfont, xspace, cancel, mathtools). The delivered numbering is the unit's own. Assets: no retained span contains a figure environment, a graphics-inclusion command, a PDF-page inclusion, a table or a TikZ picture; no image file is copied into this package, the package declares no asset dependency and no image file is redistributed with it. Licence: version arXiv:2603.28822v2 of this article is distributed under the Creative Commons Attribution 4.0 International licence, as recorded in the arXiv licence metadata for this exact version and shown on the abstract page https://arxiv.org/abs/2603.28822v2. A full rights notice is delivered at licenses/arxiv-2603.28822-CC-BY-4.0.txt. Characters: every retained excerpt is pure ASCII, so the delivered engine, XeLaTeX with its default Latin Modern text font, reports no missing character.
\begin{equation}\label{eq:zHrelation}
z_H
=
a_1+a_2+\overline{a_1a_2}\lambda,
\qquad
\lambda\in\Lambda.
\end{equation}

\begin{theorem}\label{thm:midptdist}
Let $M_1,M_2,M_3$ denote the midpoints of the sides of a triangle in $\mathcal P$, and let $O_{\mathcal D}$ be the center of the conic. Then the quantities
\[
|O_{\mathcal D}A|^2+|O_{\mathcal D}B|^2+|O_{\mathcal D}C|^2, \qquad |O_{\mathcal D}M_1|^2+|O_{\mathcal D}M_2|^2+|O_{\mathcal D}M_3|^2
\]
satisfy the orthocenter criterion.

Moreover, if the circumcenter of the triangle coincides with one of the foci of the conic, then
\[
|O_{\mathcal D}M_1|
=
|O_{\mathcal D}M_2|
=
|O_{\mathcal D}M_3|
=
\frac{R}{2},
\]
where $R$ is the circumradius of the triangle.
\end{theorem}

\begin{proof}
Assume that $\mathbb T$ is the common circumcircle of triangles in $\mathcal P$ and $a_1,a_2 \in \mathbb C$ be the foci of the conic. Then
\[
O_{\mathcal D}=\frac{a_1+a_2}{2}.
\]
Hence
\begin{align*}
|O_{\mathcal D}A|^2+|O_{\mathcal D}B|^2+|O_{\mathcal D}C|^2
&=\left|
z_1-\frac{a_1+a_2}{2}
\right|^2+\left|
z_2-\frac{a_1+a_2}{2}
\right|^2+\left|
z_3-\frac{a_1+a_2}{2}
\right|^2 \\
&=
3-\operatorname{Re}((a_1+a_2)z_H)+\frac34
|a_1+a_2|^2.
\end{align*}
Using \eqref{eq:zHrelation} gives
\[
\operatorname{Re}((a_1+a_2)z_H)=\operatorname{Re}((a_1+a_2)^2)+\operatorname{Re}(\overline{a_1a_2}(a_1+a_2)\lambda).
\]
Since $\lambda$ ranges over the admissible parameter set $\Lambda \subset \mathbb T$, the last expression is independent of $\lambda$ if and only if
\[
\overline{a_1a_2}(a_1+a_2)=0,
\]
that is, if and only if
\[
a_1+a_2=0
\qquad\text{or}\qquad
a_1a_2=0.
\]
This proves that the first quantity satisfies the orthocenter criterion.

Similarly, if $M_1$ be the midpoint of the opposite side of the vertex $z_1$. Then
\begin{align}
|O_{\mathcal D}M_1|^2
&:=\left|
\frac{z_2+z_3}{2}-\frac{a_1+a_2}{2}
\right|^2 \nonumber \\
&=
\frac14\left|z_H-z_1-(a_1+a_2)\right|^2  \nonumber\\
&=
\frac14\left|z_1-\overline{a_1a_2} \lambda\right|^2 \label{eq:odm}.
\end{align}
Summing \eqref{eq:odm} cyclically gives
\[
|O_{\mathcal D}M_1|^2+|O_{\mathcal D}M_2|^2+|O_{\mathcal D}M_3|^2
=
\frac34+\frac34|a_1a_2|^2
-\frac12
\operatorname{Re}
\left(
a_1a_2z_H\overline{\lambda}
\right).
\]
Using \eqref{eq:zHrelation} again we obtain
\[
a_1a_2z_H\overline{\lambda}
=
|a_1a_2|^2
+
a_1a_2(a_1+a_2)\overline{\lambda}.
\]
Hence
\[
|O_{\mathcal D}M_1|^2+|O_{\mathcal D}M_2|^2+|O_{\mathcal D}M_3|^2
=
\frac34+\frac14|a_1a_2|^2
-\frac12
\operatorname{Re}
\left(
a_1a_2(a_1+a_2)\overline{\lambda}
\right).
\]
The last expression is independent of $\lambda$ if and only if
\[
a_1+a_2=0
\qquad\text{or}\qquad
a_1a_2=0.
\] 
Thus the second quantity satisfies the orthocenter criterion.

Finally, if $a_2=0$, then \eqref{eq:odm} gives
\[
|O_{\mathcal D}M_1|
=
\frac12,
\]
and similarly
\[
|O_{\mathcal D}M_1|=|O_{\mathcal D}M_1|=\frac12.
\]
Thus every midpoint is at distance $1/2$ from the center of the conic. Restoring the circumradius $R$ gives
\[
|O_{\mathcal D}M_k|=\frac{R}{2},
\qquad k=1,2,3,
\]
which completes the proof.
\end{proof}

\end{document}
